% Problem 4 — Spectral Graph Theory / Free Probability (Nikhil Srivastava)
% Source: First_Proof.tex, https://github.com/1stproof/batch-1
% Shared preamble for all problems from First_Proof.tex
% Source: https://raw.githubusercontent.com/1stproof/batch-1/refs/heads/main/First_Proof.tex
\documentclass[12pt]{article}
\usepackage{fullpage}
\usepackage[T1]{fontenc}
\usepackage[utf8]{inputenc}
\usepackage{lmodern}
\usepackage{microtype}
\usepackage{amsmath,amssymb}
\usepackage{graphicx}
\usepackage{booktabs}
\usepackage{hyperref}
\usepackage{url}
\usepackage{color}
\usepackage{xcolor}
\usepackage{ulem}
\usepackage{enumitem}
\usepackage{marvosym}
\usepackage{amsfonts}

\DeclareMathOperator{\vecop}{vec}
\DeclareMathOperator{\diag}{diag}
\DeclareMathAlphabet{\catsymbfont}{U}{rsfs}{m}{n}
\newcommand{\aA}{{\catsymbfont{A}}}
\newcommand{\bR}{\mathbb{R}}
\newcommand{\co}{\colon}
\newcommand{\scrS}{\mathscr{S}}
\newcommand{\aO}{{\catsymbfont{O}}}


\title{Problem 4: Subharmonicity of $1/\Phi_n$ under Finite Free Convolution}
\author{Nikhil Srivastava (University of California, Berkeley)}
\date{}

\begin{document}
\maketitle

Let $p(x)$ and $q(x)$ be two monic polynomials of degree $n$:
\[
p(x) = \sum_{k=0}^n a_k x^{n-k} \quad \text{and} \quad q(x) = \sum_{k=0}^n b_k x^{n-k}
\]
where $a_0 = b_0 = 1$. Define $p \boxplus_n q(x)$ to be the polynomial
\[
(p \boxplus_n q)(x) = \sum_{k=0}^n c_k x^{n-k}
\]
where the coefficients $c_k$ are given by the formula:
\[
c_k = \sum_{i+j=k} \frac{(n-i)! (n-j)!}{n! (n-k)!} a_i b_j
\]
for $k = 0, 1, \dots, n$.

For a monic polynomial $p(x)=\prod_{i\le n}(x- \lambda_i)$, define
$$\Phi_n(p):=\sum_{i\le n}\left(\sum_{j\neq i} \frac{1}{\lambda_i-\lambda_j}\right)^2$$
and $\Phi_n(p):=\infty$ if $p$ has a multiple root.

The finite free convolution induces a map $$\roots: \bR^{n}\times \bR^n \to \bR^n$$, where if $\alpha$ and $\beta$ are vectors of roots for $p(x)$ and $q(x)$, then $\roots(\alpha, \beta)$ is defined to be the ordered vector of roots for $p(x)\boxplus_n q(x)$.
Given a vector of roots $\alpha\in \bR^n$ we will define the its finite free score vector $\score(\alpha) \in (\bR\cup\{\infty\})^n$ as 
$$\score(\alpha) := \left( \sum_{j : j \neq i} \frac{1}{\alpha_i-\alpha_j} \right)_{i=1}^n. $$ Given a real-rooted polynomial $p(x)$ with vector of roots $\alpha$, define its finite free Fisher information as 
$$\Phi_n(p) := \|\score(\alpha)\|^2.$$  Let $\epsilon>0$ and define the differential operator $T_\epsilon:=(1-\epsilon \cdot d/dx)^n$.

Is it true that if $p(x)$ and $q(x)$ are monic real-rooted polynomials of degree $n$, then
$$\frac{1}{\Phi_n(p\boxplus_n q)} \ge \frac{1}{\Phi_n(p)}+\frac{1}{\Phi_n(q)}?$$

\end{document}
